\scanpage{44}
define homomorphisms
\[
\tag{40}
\diagram{\begin{tikzcd}[column sep=large]
\underline M_{S_0}\arrow[r,"u_{\underline M}"] & \unsure{G^{\sharp}|G^{\natural}}
\end{tikzcd}}
,\quad
\diagram{\begin{tikzcd}[column sep=large]
\unsure{C_{S_0}|G_{S_0}}\arrow[r,"u_{\unsure{C|G}}"] & \unsure{G^{\sharp}|G^{\natural}}
\end{tikzcd}}
\quad\text{i.e.}\quad
\diagram{\begin{tikzcd}[column sep=large]
\unsure{C|G}\arrow[r] & \unsure{G^{\sharp}(S_0)|G^{\natural}(S_0)}
\end{tikzcd}}
\]
satisfying the following two conditions:

a) Commutativity of the diagram
\[
\tag{41}
\diagram{\begin{tikzcd}[column sep=large,row sep=large]
M_{S_0}\arrow[r]\arrow[d] & \unsure{C^+_{S_0}|G^+_{S_0}}\arrow[ddr,"u_{\unsure{C|G}}"] & \\
\underline M_{S_0}\arrow[drr,"u_{\underline M}"'] & & \\
& & G
\end{tikzcd}}
\quad,
\]
and b)
\[
\tag{42}
\begin{gathered}
u_{\underline M}\text{ is compatible with the actions}\\
\text{of }\unsure{C^+_{S_0}|G^+_{S_0}}\text{ on }\underline M_{S_0}
\text{ (via its action on }M_{S_0}\text{)}\\
\text{and on }G\text{ (via }u_{\unsure{C|G}}\text{).}
\end{gathered}
\]
In the present case, we will have $G=\unsure{\underline C|\underline G}$
\[
\tag{43}
u_{\underline M}(\lambda)=\deleted{\unsure{\lambda}}\,i(\mu.\lambda)
\]
where
\[
i:\quad
\diagram{\begin{tikzcd}[column sep=large]
\underline M\arrow[r,hook] & \unsure{\underline C^+|\underline G^+}
\end{tikzcd}}
\]
is the inclusion, and
\[
\tag{44}
u_{\unsure{C|G}}:\deleted{\text{defined by}}
\inserted{\diagram{\begin{tikzcd}[column sep=large]
\unsure{C^+_{S_0}|G^+_{S_0}}\arrow[r]&G
\end{tikzcd}}}
\qquad\text{defined by (33),}
\]
We must therefore first check that (33), does define a homomorphism
\diagram{\begin{tikzcd}[column sep=small]\unsure{C_{S_0}|G_{S_0}}\arrow[r]&G\end{tikzcd}},
or what amounts to the same,
\diagram{\begin{tikzcd}[column sep=small]\unsure{C|G}\arrow[r]&G(S_0)\end{tikzcd}}.
\margin{43}
\margin{43\\\texttt{https://grothendieck.umontpellier.fr}\\Université de Montpellier, UMR CNRS 5149, Case courrier 051, Place Eugène Bataillon\\34095 Montpellier cedex 5 - France}

\scanpage{45}
\mspage{\textcircled{670}}
For this we \deleted{\unsure{note that} $u$ defines a homomorphism \unsure{to the} \unsure{of index} of $3$
$\unsure{C|G}$}
\inserted{use the presentation \deleted{\unsure{of}} of $\unsure{C^+|G^+}$}
\[
\deleted{\begin{gathered}
\unsure{C_0^+|G_0^+}\\
\unsure{\simeq}\\
\unsure{\operatorname{SL}(2,\mathbb Z)}=
\{\rho,\sigma\mid\rho^6=\sigma^4=1\}
\end{gathered}}
\]
\[
\tag{44}
\unsure{C^+|G^+}=
\left\{\rho,\sigma\ \middle|\ 
\begin{gathered}
\rho^6=\sigma^4=1,\quad\rho^3=\sigma^2\\
[\sigma,\lambda_\infty]=1\quad\text{where}\quad
\lambda_\infty=\deleted{\unsure{\rho}}\rho\sigma\rho^{-1}\sigma\rho
\end{gathered}
\right\},
\]
thus setting
\[
\tag{45}
\left\{\begin{aligned}
\rho'&=\unsure{\zeta\rho|u(\rho)},\quad
\sigma'=\eta_0\omega^{\nu_2}\sigma\qquad\text{cf. (33) and}\\
\unsure{\zeta|\xi}&=\alpha_0-\rho(\alpha_0)+\nu_3\lambda_1,\quad
\unsure{\eta|\eta_0}=\beta_0-\sigma(\beta_0)
\end{aligned}\right.
\]
we must check the relations
\[
\tag{46}
\rho'{}^6=\sigma'{}^4=[\sigma',\lambda'_\infty]=1,\quad
\rho'{}^3=\sigma'{}^2
\quad\text{where}\quad
\lambda'_\infty=\deleted{\unsure{\rho'(\lambda'_0)}}
\inserted{\rho'\sigma'\rho'{}^{-1}\sigma'\rho'}.
\]
We will begin by checking \deleted{\unsure{that|more}}
\[
\tag{47}\rho'{}^3=\sigma'{}^2=\omega
\]
which implies \deleted{\unsure{that}}
\[
\tag{48}
\rho'{}^6=\sigma'{}^4=1,\quad\text{thus}\quad
\exists\ \deleted{\text{unique}}
\diagram{\begin{tikzcd}[column sep=large]
\unsure{C_0^+|G_0^+}\arrow[r,"u_{\unsure{C_0|G_0}}"] & G(S_0)\\
\rho\arrow[r,mapsto]&\rho',\quad\sigma\mapsto\sigma'
\end{tikzcd}}
\]
Next, \deleted{setting} noting that we will have
\[
\lambda'_\infty=u_{\unsure{C_0|G_0}}(\lambda_\infty),
\]
it will suffice to \deleted{\unsure{this}} check
\[
\tag{49}\lambda'_\infty=\mu\lambda_\infty
\]
it will be obvious that $\lambda'_\infty$ commutes with
\[
\sigma'=u(\sigma)=\eta_0\sigma\qquad\deleted{\text{\unsure{pui}}}
\]
since $\sigma$ and $\eta_0\in M$ commute with it. Thus to construct
$u_{\unsure{C|G}}$, the essential checks are (47) and (49)
--- moreover once (47) is proved, the preceding calculations leading to (38)
are justified (\unsure{at least} where they make sense, i.e.,
for $u_{\unsure{C|G}}$ and $\lambda\in M$),
\margin{44}
\margin{44\\\texttt{https://grothendieck.umontpellier.fr}\\Université de Montpellier, UMR CNRS 5149, Case courrier 051, Place Eugène Bataillon\\34095 Montpellier cedex 5 - France}

\scanpage{46}
so (49) will indeed be verified.

We \deleted{will have} have
\[
\begin{aligned}
&\deleted{\unsure{u(\rho)=}}\\
\rho'{}^3&=\bigl(\unsure{\zeta|\xi}\rho(\unsure{\zeta|\xi})\rho^2(\unsure{\zeta|\xi})\bigr)\rho^3
=\unsure{\zeta|\xi}\rho(\unsure{\zeta|\xi})\rho^2(\unsure{\zeta|\xi})\,\omega\\
\sigma'{}^2&=(\eta_0\sigma(\eta_0))\sigma^2=(\eta_0\sigma(\eta_0))\omega
\end{aligned}
\]
so \deleted{and} formulas (47) reduce to:
\[
\tag{50}
\unsure{\zeta|\xi}\rho(\unsure{\zeta|\xi})\rho^2(\unsure{\zeta|\xi})=1,
\qquad\eta_0\sigma(\eta_0)=1
\]
i.e., additively
\[
\unsure{\zeta|\xi}+\rho(\unsure{\zeta|\xi})+\rho^2(\unsure{\zeta|\xi})=0,
\qquad\eta_0+\sigma(\eta_0)=0.
\]
whose verification is trivial, given that
\[
\rho_{\underline M}^3=\sigma_{\underline M}^2=\mathrm{id}_{\underline M},
\qquad(1+\rho_{\underline M}+\rho_{\underline M}^2)
=\deleted{\unsure{\mathrm{id}_{\underline M}}}0_{\underline M}.
\]
We have thus indeed constructed a homomorphism
\deleted{\unsure{to the}}
\[
u_{\unsure{C_0^+|G_0^+}}:\quad
\diagram{\begin{tikzcd}[column sep=large]
\unsure{C_0^+|G_0^+}\arrow[r]&\unsure{\underline C_0^+(S_0)|\underline G_0^+(S_0)}
\end{tikzcd}},
\]
and it remains to check (41) and (42).

For (41), we have already shown it at the \inserted{level of the existence of a $u_{\unsure{C_0|G_0}}$}
--- it amounts to proving
\[
\bigl(\quad u_{\unsure{C_0|G_0}}(\lambda_0)=\deleted{\unsure{\lambda_0}}\mu\lambda_0,
\quad u_{\unsure{C_0|G_0}}(\lambda_1)=\mu\lambda_1
\]
I will all the same redo the verification, to be sure.
\[
\begin{aligned}
\varepsilon'_0&=u_{\unsure{C_0|G_0}}(\underbrace{\varepsilon_0}_{\sigma\rho})
=u_{\unsure{C_0|G_0}}(\sigma)u_{\unsure{C_0|G_0}}(\rho)
=\eta_0\omega^{\nu_2}\deleted{\unsure{\zeta|\xi}}\sigma\unsure{\zeta|\xi}\rho
=\eta_0\omega^{\nu_2}\sigma(\unsure{\zeta|\xi})\sigma\rho\\
&=\eta_0\sigma(\unsure{\zeta|\xi})\omega^{\nu_2}\varepsilon_0
\end{aligned}
\]
\margin{45}
\margin{45\\\texttt{https://grothendieck.umontpellier.fr}\\Université de Montpellier, UMR CNRS 5149, Case courrier 051, Place Eugène Bataillon\\34095 Montpellier cedex 5 - France}

\scanpage{47}
\mspage{\unsure{671}}

I claim that
\[
\begin{aligned}
\varepsilon'_0&=\varepsilon_0^{(\mu)}\quad\text{i.e.}\quad
\eta_0\sigma(\unsure{\zeta|\xi})\omega^{(\nu_2)}
 =\underbrace{\varepsilon_0^{(\mu-1)}}_{l_0^\nu}
\quad\text{where }\mu-1=2\nu\deleted{\unsure{+1|=1}},\quad\text{i.e.}\\
\eta_0\sigma(\unsure{\zeta|\xi})\omega^{\nu_2}&=l_0^{(\nu)}
\quad\text{where }\nu=2\nu'\unsure{-|+}\nu_2\quad\text{i.e.}\\
\eta_0\sigma(\unsure{\zeta|\xi})
 &=\underbrace{l_0^{(2\nu')}}_{\lambda_0^{2\nu'}}
 \underbrace{\unsure{(\omega l_0)^{(-\nu_2)}|(\omega l_0^{(-\nu_2)})}}_{\lambda_0^{-\nu_2}}
 =\underbrace{\lambda_0^{(2\nu'\unsure{-|+}\nu_2)}}_{\lambda_0^\nu}
 \quad\text{i.e.}\\
\eta_0\sigma(\unsure{\zeta|\xi})&=\lambda_0^{(\nu)}
\end{aligned}
\]
i.e., applying $\sigma$ and noting that
$\sigma(\eta_0)=\eta_0^{-1}$, $\sigma^2(\unsure{\zeta|\xi})=\unsure{\zeta|\xi}$,
$\sigma(\lambda_0)=\lambda_1$,
\[
\eta_0^{-1}\unsure{\zeta|\xi}=\lambda_1(\nu)
\quad\text{i.e. finally }\deleted{\unsure{\zeta_0=|\xi_0=}}
\unsure{\zeta|\xi}=\eta_0\unsure{\lambda_1(\nu)|\lambda_1^{(\nu)}}
\quad\text{or}
\]
\[
\unsure{\zeta|\xi}=\eta_0+\nu\lambda_1\quad\text{(written additively)}
\]
which is true by hypothesis (cf. (31)).

We will then have
(\deleted{\unsure{for $u=u_{\unsure{C_0|G_0}}$|for $u_{\unsure{C_0|G_0}}$}}
\inserted{writing $u$ instead of $u_{\unsure{C_0|G_0}}$ to simplify})
\[
u_{\deleted{\unsure{C_0|G_0}}}(\lambda_1)
 =u\rho(\lambda_0)
 =(\unsure{\zeta|\xi}\rho u)(\lambda_0)
 =\unsure{\zeta|\xi}\rho(u(\lambda_0))
 =\unsure{\zeta|\xi}\rho(\lambda_0(\mu))
 =\lambda_1(\mu)
\]
OK.

It remains to check (42), which \deleted{\unsure{means}}
\inserted{amounts} to the \unsure{commutativity} of
$u_{\underline M}=\mu\,\mathrm{id}_{\underline M}$ with
$(\rho_{\underline M},\rho'_{\underline M})$,
$(\sigma_{\underline M},\sigma'_{\underline M})$
--- \unsure{this means} also that
$\rho_{\underline M}$ and $\rho'_{\underline M}$ on the one hand,
$\sigma_{\underline M}$ and $\sigma'_{\underline M}$ on the other hand,
\deleted{\illegible{2}}\inserted{are equal, which}
is also obvious.

Thus, the group $\widetilde Z_{\rho,\sigma}$ is well and
truly constructed --- at least, as a
scheme over $S_0=\mathrm{Spec}\,\mathbb Z$, with morphisms
of schemes
\[
\tag{51}
\diagram{$
\begin{tikzcd}[column sep=3em,row sep=.25em]
\widetilde{\underline Z}_{\rho,\sigma}\arrow[r,"\nu'"] & G_a\\
{}\arrow[r,"\deleted{\unsure{\nu|\varepsilon}}"] & \deleted{\unsure{\{\pm1\}_{S_0}^{2}|\mu_{S_0}^{2}}}\\
{}\arrow[r,"\nu_2"] & \{0,1\}_{S_0}\\
{}\arrow[r,"\nu_3"] & \{0,1\}_{S_0}
\end{tikzcd}
\quad\Bigg|\quad
\begin{gathered}
\text{whence }\mu=4\nu'+2+\nu_2:\quad
\begin{tikzcd}[column sep=1.5em,row sep=1em]
\widetilde Z_{\rho,\sigma}\arrow[r]\arrow[rd] & G_a\\
 & G_m\arrow[u,hook]
\end{tikzcd}\\
\varepsilon_2,\varepsilon_3:\quad
\begin{tikzcd}[column sep=3em]
\widetilde Z_{\rho,\sigma}\arrow[r,shift left=.3ex]\arrow[r,shift right=.3ex]
 & \mu_{S_0}=\{\pm1\}_{S_0}
\end{tikzcd}
\end{gathered}
$}
\]

\margin{46}
\margin{46\\\texttt{https://grothendieck.umontpellier.fr}\\Université de Montpellier, UMR CNRS 5149, Case courrier 051, Place Eugène Bataillon\\34095 Montpellier cedex 5 - France}